基于对1-Def.v1_polished_unified_expanded.tex文件的分析，以下是我诊断出的需要扩写1-1.5倍的段落及其扩写版本：

【原文1】
当直觉被形式化后，我们究竟"保留下了什么"，又"舍弃了什么"？现实生活中一个很直观的例子，是我们在疫情期间依然能轻易认出戴着口罩的熟人：即便只露出眼睛和部分脸部轮廓，我们往往仍能在极短时间内判断"这是谁"。

【扩写版本1】
当直觉被形式化转化为数学算法时，一个值得深思的问题是：我们究竟"保留下了什么"，又"舍弃了什么"？这个问题的答案揭示了人类智能与机器智能之间的本质差异。

让我们来看一个生活中非常直观的例子。在疫情期间，我们依然能够轻易认出戴着口罩的熟人。尽管对方的面部大部分被遮挡，只露出眼睛和部分脸轮廓，但我们往往仍能在极短时间内准确判断出"这是谁"。这种能力为什么如此神奇？因为人类的大脑并不仅仅依赖于完整的面部特征，而是会综合运用多种线索：眼神的独特性、眉宇间的神情、脸型的轮廓，甚至是对方走路姿态或说话语气的细微特征。

然而，当人工智能试图将"识别人脸"这种人类直觉转化为算法时，这个过程就变得截然不同。计算机会怎么做呢？它通常会将图像转化为像素矩阵，通过数学运算提取数值特征，然后与数据库中的已知面孔进行匹配。这个过程中保留下的是像素之间的数值关系与统计模式，舍弃的则是人类在注视他人时所携带的丰富情感背景、社会语境和个人记忆。

【原文2】
正是数学的引入——尤其是统计学、优化理论与概率思想——让机器开始具备了从数据中学习、在经验中改进的能力。通过在大量样本中寻找规律、量化误差并不断优化，智能真正跨入了"学习时代"。

【扩写版本2】
正是数学的引入——尤其是统计学、优化理论与概率思想的系统性融合——让机器开始具备了从数据中学习、在经验中改进的能力。这个转变为什么如此重要？因为它标志着AI从"死板执行"转向"智能适应"的根本性飞跃。

让我们具体理解这个过程。统计学教会机器如何从大量样本中发现隐藏的规律和模式；优化理论提供了不断改进性能的方法论；概率思想则让机器能够在不确定性中做出合理的推断。通过在这些数学工具的指导下，机器能够在海量样本中寻找规律、量化预测误差并不断优化自身参数，智能真正跨入了"学习时代"。

这就像一个学生从只会背诵标准答案，到真正理解了知识并能够举一反三的转变。早期的人工智能系统就像是只会背诵答案的学生——它们严格按照预设规则执行任务，一旦遇到新情况就无法应对。而有了数学学习能力的现代AI系统，则像是真正理解了知识的学生——它们能够从经验中总结规律，在新的情境中灵活应用，并且随着经验积累而不断进步。

【原文3】
鹦鹉是自然界的"模仿大师"。它能精准复述人类语言，甚至在特定场景中给出看似"理解"的回应。然而，它并不真正懂得自己在说什么——它模仿的是声音的模式，而无法将词语与真实世界中的对象建立语义联系。

【扩写版本3】
鹦鹉堪称自然界的"模仿大师"，它的语言模仿能力令人惊叹。它能精准复述人类的话语，甚至在某些特定场景中给出看似"理解"的回应，让人误以为它真的在思考。然而，鹦鹉行为的本质究竟是什么？它并不真正懂得自己在说什么，而是在模仿声音的声学模式。

这个现象为什么值得深入思考？因为它揭示了"表面智能"与"真正理解"之间的关键区别。鹦鹉无法将听到的词语与真实世界中的具体对象建立语义联系。当它说"苹果"时，它只是在模仿这个词汇的发音，而脑海中并没有苹果的形象、味道或营养价值的概念。

这种"鹦鹉式智能"的特征与当前许多人工智能系统惊人地相似。特别是大型语言模型，它们能够生成语法正确、语义连贯的文本，甚至进行复杂的对话，表面上看起来像是真的在理解。但就像鹦鹉一样，这些模型处理的是符号的统计模式，而非概念的真实含义。它们通过计算概率来选择下一个最可能出现的词汇，而非基于对世界的真正理解来生成回应。

【原文4】
从计算复杂度角度看，二者的差异也十分明显。鹦鹉式智能通常需要大量的计算资源与训练数据，但一旦训练完成，推理过程相对直接。相反，乌鸦式智能在学习阶段可能只需少量数据，却在推理时承担更高的复杂度，需要动态的因果推理和创造性思维。

【扩写版本4】
从计算复杂度的角度来看，这两种智能模式的差异也十分明显，几乎呈现出完全相反的计算模式特征。为什么这样说呢？因为它们在学习和推理阶段的资源分配策略截然不同。

鹦鹉式智能通常需要投入大量的计算资源与训练数据才能达到较好的性能。在训练阶段，它必须处理数以亿计的样本，通过反复的试错和参数调整来掌握任务的统计规律。这个过程计算量巨大，耗时极长，而且对硬件要求很高。然而，一旦训练完成，在实际推理时，系统只需要进行相对直接的前向计算即可得出结果，推理过程较为高效和稳定。

与此形成鲜明对比的是，乌鸦式智能在学习阶段往往只需少量样本就能快速掌握问题本质。就像乌鸦通过几次观察就能学会使用工具一样，真正的智能体具备快速学习和举一反三的能力。但是，在进行实际推理时，它需要承担高得多的计算复杂度，因为它必须进行动态的因果推理、创造性思维和灵活的问题解决。

这种差异反映了两种不同的智能哲学：鹦鹉式智能追求"训练复杂，推理简单"，而乌鸦式智能则体现"训练简单，推理复杂"的特点。这就像两种不同的学习方法——一种是题海战术，通过大量练习来掌握固定模式；另一种是理解性学习，通过掌握基本原理来解决各种新问题。
